\section{Discussion and Limitations}\label{sec:discussion}
\subsection{Interpreting Reference Alignment}
Common-model operating regimes and descriptive native-output similarity address RQ1; certificates and forward/inverse feasibility queries address RQ2. The tested reference-order gains extend across planar and volumetric inputs, while motion and geometry gains depend on the reference and event population (RQ3). Nucleon's unchanged P/R profiles show that useful alignment need not require relaxation.

A band's movement combines projected source movement and changes in its anchor-to-reference residual. Wildfires frames 17--19 show a fixed reference accompanying a shifting P anchor; matched source views separate these effects. For volumetric fields, the reference compresses two further spatial dimensions, so similar projected positions can represent features far apart in the source domain.

The new common-model experiment compares design principles directly within one representation. It also reveals their limits: native-output resemblance varies by input, with substantially weaker Storms ST-MTM agreement than Ring. The anchor-level temporal surrogate omits all-subtree overlap and uses shared absolute widths, so it should not be interpreted as an implementation substitute for either native method.

Reference alignment complements distance preservation and temporal coherence. On OLR, P/R motion-proxy errors remain 3/56 despite improved order, and ST-MTM retains better SNS, TW, and TD on Storms, OLR, and PM2.5. Annual PM2.5 has 95 clear events at the primary thresholds, with 42 P errors and 47 R20 errors. These outcomes preclude treating a lower NROE as a general guarantee of better geometry, motion interpretation, or reader performance. The motion-margin result in Supplement S1 instead certifies projected distance changes when achieved anchor errors are sufficiently small.

\subsection{Choosing Preferences and Error Budgets}
A user selects the relations to emphasize within the common interval model. Setting reference and geometry weights to zero, disabling the position cap and using a stationary temporal target gives topology--time alignment. Enabling centroid stress gives topology--geometry preservation. Reference guidance activates absolute targets and projected motion; a scalar allowance determines how much merge-level change can be spent to resolve order conflict. These regimes share one construction. Native TMTM and ST-MTM remain independent baselines with different sample allocation and width rules, so a shared design principle does not establish output identity. An application can prescribe either error cap or ask the inverse query for the allowance required by a position tolerance.

RG-MTM-P is appropriate when source merge thresholds must remain intact. Its position certificate quantifies the displacement that the chosen geometry and hierarchy permit, while its actual anchor errors disclose the result of the budgeted secondary layout. RG-MTM-R adds reference-compatible orders while retaining the extrema and their values. Each relaxed output should be read together with its achieved displacement and mean and maximum merge-level errors; sequence means alone can conceal a large local scalar change.

The eligibility parameter $\kappa$ bounds adjacent parent gaps of nodes that may be flattened. It is neither birth--death persistence nor an output-distortion tolerance. Wider eligibility need not improve the measured trade-off: Cylinder R100 has the same mean NROE as R20 but greater mean NMLE. The new RG-MTM-B specialization imposes that tolerance on continuous fixed-order costs and rejects inadmissible candidates. Its 603 exhaustive frames supply exact optima; 11 candidate-only frames supply feasible upper bounds. Strict-extremum raster output still needs explicit clearance and actual-error acceptance checks. The retained large-tree eligibility construction supplies separate temporal maps, not the new hard-budget guarantee.

Relaxation changes regular scalar values and connection thresholds. The negative Wildfires connections are artificial filling values, even though the source fire observations are nonnegative. Merge-level error describes these changed pairwise thresholds, while retaining extremum values; it does not establish preservation of scalar distributions, thresholded areas, or every internal feature. An analysis that depends on a physical threshold should consult the source field and the displayed barrier before drawing a conclusion from the relaxed map.

\subsection{Scope of the Evidence}
Spatial projection and preprocessing determine the structures represented by the map. Tangaroa's direction captures 57.686\% of matched motion energy; exact projected anchors still produce 2060/6755 three-dimensional motion-proxy errors. Cylinder's exact $x$ projection likewise retains 603/3196 errors. These controls locate a limit of the chosen reference even when positional layout error vanishes. Multiple references could expose additional spatial relations, but are outside the present evaluation.

New budget experiments use the same 614 planar frames but construct new LP/barrier witnesses; they do not reuse previous temporal scores as measurements of those witnesses. The reported populations answer different questions. The 614-frame suite compares four methods on common planar temporal inputs. The 64-frame mirror-motion control isolates absolute anchoring. Tangaroa and Cylinder test independent original-grid volumetric inputs; Nucleon tests one static slice; reconstructed Benzene tests section coherence in a declared classical potential. The 744-hour Storms overview supplies a denser visualization and is excluded from the 118-frame quantitative sequence. These populations are not combined into an overall accuracy score. ST-MTM was not run on the static additions or volumetric sequences.

Input correspondence is imperfect. Tangaroa has 2459/6925 boundary-touching supports and 36--186 derived nodes (reference: 50--186). Cylinder uses 102 public volumes rather than the reported 508 frames. Benzene is a reconstructed potential; annual PM2.5 replaces an exploratory January sample. These declared protocols support the reported comparisons, not exact reproduction of every original input or untested generalization.

Overlap matches provide computational correspondences rather than validated physical tracks. Source-provider event calendars give temporal context rather than cyclone identities, and the Cylinder proxies do not establish vortex labels. Independent physical-track validation and studies of how readers interpret the maps remain necessary extensions. Further persistence settings, internal-feature evaluation, additional scalar models, and preregistered task-specific comparisons would broaden the evidence.

\subsection{Computational Limits and Future Work}
Fixed-order position costs and non-strict barriers are exact within the model; complete enumeration yields exact frontiers. The common solver evaluates every admitted candidate, while the retained large-tree construction uses brackets and at most four QPs. Both can leave search gaps. Finer previews do not strengthen numerical evidence.

Future work includes richer overlap-aware temporal terms, scalable search, sequence-wide optimization and verified raster acceptance. Historical planar QPs may fall back to LP-feasible layouts; their optimality and end-to-end timing need further assessment. The ST-MTM SciPy OLO heuristic exceeds the exhaustive adjacency minimum in eight of 148 audited frames, qualifying the declared reimplementation rather than every original setting.
